The following terms are central to the conceptual and methodological framework of this dissertation. Some are established in existing scholarship, while others are introduced here as neologisms to name processes and relationships that existing terminology does not adequately capture. Where relevant, terms are linked to prior scholarship while also clarified in terms of their specific usage in this project.

\subsection*{Core Framework Terms}

\begin{description}
	\item[Embodying]: A phase of the Speculocultural Technopoietic cycle concerned with testing and validating modified toolsets through practice and creation. Embodying prioritizes cultural authenticity over technical correctness, asking whether modifications feel genuine to the cultural knowledge that guided their creation and whether they enable forms of expression that honor their epistemological foundations.
	
	\item[Fugitive Speculation] (\textit{neologism}): A mode of speculative thought rooted in refusal, improvisation, and resistance. It emphasizes imagining otherwise within systems of constraint, drawing from traditions of fugitivity to unsettle dominant logics. \textit{See also}: Fugitive Epistemologies.
	
	\item[Intimate Witnessing] (\textit{neologism}): A form of personal and collective witnessing that validates cultural knowledge, acknowledges vulnerability, and emphasizes ethical relations within collaborative practice.
	
	\item[Material Memory] (\textit{neologism}): The recognition that objects, media, and technologies act as cultural archives, carrying accumulated histories and resonances of the past into the present. Material memory describes the archival capacity of artifacts themselves, which can be mobilized through countermemory practices.
	
	\item[Modification]: The deliberate, culturally-informed practice of altering the technical, algorithmic, and interactional structures of tools to embed alternative epistemologies and resist dominant design paradigms. In the context of this research, modification operates at multiple levels: (1) code-level (rewriting algorithms, reconfiguring data structures, altering parameter mappings); (2) interface-level (redesigning affordances, interaction patterns, and representational schemas); and (3) conceptual-level (reframing what counts as valid input, output, or creative gesture).
	
	\item[Remixing]: A phase of the Speculocultural Technopoietic cycle characterized by embodied modification and material manipulation that results in knowledge creation. Remixing encompasses the development of alternatives and the discovery of unauthorized approaches, generating knowledge by making technologies "work differently" through specific cultural epistemologies.
	
	\item[Sensory Worlding] (\textit{neologism}): Ways of knowing and creating rooted in multisensory engagement. This term emphasizes embodied, non-ocularcentric perception as a method of constructing worlds and designing technologies.
	
	\item[Speculocultural] (\textit{neologism}): A conceptual term coined in this dissertation to describe the speculative reimagining of technology through specific cultural epistemologies. Like \textit{sociocultural}, it marks the conjoining of domains: speculation as imaginative resistance, and culture as epistemological grounding.
	
	\item[Speculocultural Technopoiesis] (\textit{neologism}): The integrated framework of this dissertation. It names the cyclical relationship between cultural speculation and embodied technological practice, where imagination guides making and making generates new knowledge that reshapes cultural imagination.
	
	\item[Transmitting]: A phase of the Speculocultural Technopoietic cycle where successful cultural-technical discoveries become methodologies that members of kindred communities can adapt with relative ease and understanding. Transmission preserves not only techniques but the cultural reasoning behind them.
	
	\item[Unsettling]: A phase of the Speculocultural Technopoietic cycle that arises when embodied lived experience encounters, often unsatisfactorily, the hidden assumptions of technological toolsets. Unsettling involves encountering and categorizing the frictions between cultural embodied knowledge and technological assumptions, revealing biases embedded in existing paradigms.	
\end{description}

\subsection*{Methodological Terms}

\begin{description}
	\item[Autoethnography]: A research methodology that uses the researcher's personal narrative as both process and product of analysis.
	
	\item[Collaborative Autoethnography]: A small-cohort practice where multiple participants share and reflect on their lived experiences, creating situated knowledge collectively rather than individually.
	
	\item[Cyclical Praxis]: An approach to practice that foregrounds iteration, return, and repetition as generative methods. Cyclical praxis embraces reworking and revisiting as acts of resistance and care.
	
	\item[Praxis]: The integration of theory and practice, where knowledge is generated through action and reflection. Praxis emphasizes that understanding emerges not from abstract thought alone, but through engaged doing.
	
	\item[Situated Knowledge]: Knowledge that is recognized as partial, embodied, and produced from specific social and material locations. Introduced by Donna Haraway, situated knowledge rejects claims to universal objectivity in favor of accountable, positioned perspectives.
\end{description}

\subsection*{Theoretical Foundation Terms}

\begin{description}
	\item[Afrofuturism / Black Speculative Traditions]: Aesthetic and critical frameworks that explore intersections of imagination, technology, the future, and liberation within Black cultural practices.
	
	\item[Agential Realism]: A framework developed by Karen Barad that emphasizes the entanglement of matter and meaning. It views reality as enacted through "intra-actions" rather than interactions between pre-existing entities.
	
	\item[Decoloniality]: A theoretical and practical orientation toward undoing colonial power structures, epistemologies, and ontologies. Decolonial thought emphasizes pluriversality, the legitimacy of multiple ways of knowing and being, and the dismantling of Eurocentric universalism.
	
	\item[Epistemology]: The study of knowledge itself: how we come to know, what counts as knowledge, and who is authorized to produce it. In this dissertation, epistemology is understood as culturally situated and materially grounded rather than abstract or universal.
	
	\item[Fugitive Epistemologies]: Knowledge practices that refuse dominant paradigms, emerging through improvisation, brokenness, and marginal positioning. These epistemologies arise from traditions of fugitivity and marronage, creating ways of knowing that evade capture by hegemonic systems. \textit{See also}: Fugitive Speculation.
	
	\item[Intra-action]: Karen Barad's term for the mutual constitution of entities through their entanglement. Unlike "interaction" (which presumes pre-existing separate entities), intra-action emphasizes that boundaries and properties emerge through relational encounters rather than existing beforehand.
	
	\item[Poiesis]: From the Greek, meaning "making" or "bringing forth." Poiesis refers to creative production that brings something into being, emphasizing the generative and transformative dimensions of making. The term underscores that creation is not merely technical execution but a form of knowledge production.
	
	\item[Sonic Fiction]: Introduced by Kodwo Eshun, this term describes how Black electronic music creates speculative futures, treating frequencies and rhythms as narrative devices and world-making practices.
	
	\item[Technopoiesis]: Knowledge creation through embodied technological making and material manipulation. Developed by Guerra and Ostergaard, the concept describes the cyclical co-evolution of human, material, and technical knowledge through acts of making.
	
	\item[The Undercommons]: A concept from Fred Moten and Stefano Harney describing fugitive planning, brokenness, and shared debt as principles of knowledge and community.
\end{description}

\subsection*{Practice-Based Terms}

\begin{description}
	\item[Call-and-Response]: A dialogic pattern rooted in African and African diasporic musical and oral traditions, where a statement (call) is answered by a reply (response). In this dissertation, call-and-response operates as both a sonic structure and an epistemological principle that shapes collaborative computation and cultural encoding.
	
	\item[Collaborative Computation (Sonic)]: Collective creation of algorithms and sound systems through distributed decision-making and cultural exchange.
	
	\item[Countermemory]: Memory practices that resist dominant or colonial archives, surfacing suppressed or marginalized histories. It actively contests official narratives by centering alternative witnesses, experiences, and material traces.
	
	\item[Cultural Encoding (Sonic)]: Embedding cultural sonic logics into computational systems, ensuring that algorithms reflect situated traditions of rhythm, call-and-response, or improvisation.
	
	\item[Embodied Knowledge]: Knowledge that arises through bodily practice, sensation, and movement, rather than abstract or disembodied reasoning.
	
	\item[Improvisation]: A mode of real-time creative decision-making that responds to emergent conditions. In Black musical traditions and this dissertation's framework, improvisation is not randomness but a disciplined practice drawing on deep cultural knowledge and relational listening.
	
	\item[Ocularcentrism]: The privileging of vision and visual knowledge above other sensory modalities. Ocularcentric paradigms dominate Western epistemology and design, often marginalizing sonic, haptic, and other embodied ways of knowing.
	
	\item[Remix]: A method of recombination and transformation of existing materials. Within this dissertation, remix operates as both artistic practice and epistemological strategy, generating new knowledge through creative recombination.
\end{description}